The guarantee is a certificate; whether it also saves money depends on how many calibration units it rests on. Proposition~\ref{prop:cost}(ii) gives the mechanism: a level chosen among $n$ calibration critical levels fails with probability at least $1/(n+1)$, and with expensive failures that floor dominates the cost. With the split design and 32 calibration units the certified alarm was clearly dearer than the best uncertified rules; cross-fitting, which calibrates on all training units and trains the signal on four fifths of them, removed most of the premium (Table~\ref{tab:cost}). With equal information, cross-fitted calibration made the certificate almost free: the certified alarm was never significantly dearer than the best uncertified rules built on the same models, and cheaper than the one-standard-error rule on two datasets (Section~\ref{sec:res-cost}).
 Rules that choose levels beyond the most demanding calibration unit can be cheaper on a given test set, and the earlier study showed how such rules can fail when plain empirical cost minimisation finds a setting at the edge of the feasible region~\cite{prior}; what they cannot provide is a bound for the next unit. Two sizing rules follow. To certify a failure probability $\alpha$ one needs at least $1/\alpha-1$ run-to-failure (or late-censored) calibration units from the population in question; to keep the certified excess cost $(\rho-1)/(n+1)$ below a fraction $\varepsilon$ of a planned replacement one needs $n+1\ge(\rho-1)/\varepsilon$. Where such data exist, as for fleets of identical components or cells from one production line, the certificate is cheap; where they do not, a failure rate of an uncertified alarm measured on a few dozen units is weak evidence of its reliability.
